Shelfmark 149, scan 121.

\scanpage{121}
\mspage{60}
This being so, $\mathcal E_U^{\overline L}$ is recovered from $\mathcal E_L^{\overline L}$, as a quotient of the latter, by taking \uline{all} the $V$ in $L$ having a centre on $U$ (it suffices $=$ to take the $V=\underline{O}_{U,x}$, where $x$ is a point of codimension 1 in $U$), and dividing by the closed invariant subgroup generated by the corresponding $I_{x'}$.

One can consider the locally ringed space, $\varprojlim$ of the \inserted{(regular)} proper models of $L$ over $K$
\[
(31)\qquad\mathfrak X_{L/K}\simeq\varprojlim_{\substack{X\text{ proper model}\\\text{of }L/K}}X
\]
\margin{\uline{Note well} This definition does not depend on the hypothesis of resolution of singularities --- for one can take normal models --- not necessarily regular --- (It amounts to the $=$, if one has resolution\ldots)}
--- this is \uline{not} a scheme ! \inserted{(if $n\geq2$)}. \deleted{\unsure{Locus|One}} \inserted{If}
$x\in\mathfrak X$, \deleted{\unsure{the set of|the ring of}} \inserted{$k(x)$} is an extension of $K$ of transcendence degree $d\leq n$ --- one calls it the \uline{dimension} of $x$, \inserted{and $n-d$ is the \uline{codimension}} There is one and only one point of codimension 0 --- corresponding to the generic points of the $X$. The points of codimension 1
\margin{120}
